\section{Heavy modes and the reduction to one-sided error}\label{sec:heavy}

The two discrepancy signs have different obstructions. An omitted mode
contributes its entire terminal mass to $A-W$, whereas a long activation
block can make $W-A$ large. The following theorem separates these effects
for every block count. A mode is called \emph{heavy at threshold $E$} when
$m_i>E$; the inequality is strict.

\begin{theorem}[Universal heavy-mode rounding]\label{thm:heavy}
Let $s\ge0$, $E>0$, and $T\le(s+2)E$. If some mode has $m_i>E$,
then there is a schedule with at most $s$ switches and full error at most $E$.
There is no restriction relating $n$ to $s$.
\end{theorem}

We prove the theorem by rounding prefixes on equal intervals and then
reordering a carefully chosen prefix. This reordering is needed because a
repeated mode alone does not save a switch: its occurrences must be adjacent.

\begin{lemma}[Integral prefix rounding with a repetition]\label{lem:prefix-repeat}
Let $A\in\Acal_n(M)$, where $M\ge2$ is an integer. If $A_h(M)>1$,
there is a word $w_1,\ldots,w_M$ on unit intervals whose counts $N_i(j)$ obey
\begin{equation}\label{eq:integral-prefix}
 \lfloor A_i(j)\rfloor\le N_i(j)\le\lceil A_i(j)\rceil
 \qquad(i\in[n],\ 1\le j\le M),
\end{equation}
and in which mode $h$ occurs at least twice. This word has full error at most
one for the original measurable input.
\end{lemma}
\begin{proof}
Build a finite network with one supply-one node for each slot $j$ and, for
each mode $i$, a chain of nodes $(i,1),\ldots,(i,M)$. Slot $j$ has an arc
to each $(i,j)$, with capacity one. The chain edge from $(i,j)$ to
$(i,j+1)$ has lower and upper capacities $\lfloor A_i(j)\rfloor$ and
$\lceil A_i(j)\rceil$; the last chain edge goes to a common sink with
the corresponding capacities at $M$. The sink has demand $M$. Sending
$A_i(j)-A_i(j-1)$ from slot $j$ to its mode chain gives a feasible
fractional flow. All bounds and supplies are integers. The integral-flow
theorem (equivalently, total unimodularity of a directed incidence matrix)
gives an integral optimum for a linear objective on this bounded feasible
flow polytope.

Maximize the final flow in chain $h$. Its optimum is at least the feasible
fractional value $A_h(M)>1$, so an integral optimum assigns at least two
slots to $h$. Each slot sends its one unit to exactly one mode, and its
chain counts satisfy~\eqref{eq:integral-prefix}. The endpoint error is at
most one. Lemma~\ref{lem:endpoint} extends this bound throughout each
unit interval without any assumption that the input is constant there.
\end{proof}

\begin{lemma}[First-repeat reordering]\label{lem:first-repeat}
Any word on $M$ unit intervals with full error at most one and a repeated
mode can be replaced by a word with full error at most one and an equal
adjacent pair.
\end{lemma}
\begin{proof}
If there is already an adjacent pair, nothing is required. Otherwise let
$r$ be the first position repeating an earlier mode $q$. The prefix of
length $r$ contains $q$ twice and $r-2$ other modes once each. Its error
at $r$ gives
\begin{equation}\label{eq:repeat-prefix-masses}
 1\le A_q(r)\le3,\qquad
 A_i(r)\le2\ \text{for other selected modes},\qquad
 A_i(r)\le1\ \text{for omitted modes}.
\end{equation}
Only this prefix will be reordered, preserving its occupation multiset.

Let $d_q=\max\{t\in[0,r]:A_q(t)\le1\}$. Continuity and
\eqref{eq:repeat-prefix-masses} give $A_q(d_q)=1$, including the case
$A_q(r)=1$ where $d_q=r$. Moreover $1\le d_q\le r$, since $A_q(t)\le t$.
Choose an integer $j\in\{0,\ldots,r-2\}$ with $j\le d_q\le j+2$;
one choice is $j=\min\{r-2,\max\{0,\lceil d_q\rceil-2\}\}$.
Place $q$ on the double block $[j,j+2]$.

For another selected mode $i$ with $A_i(r)>1$, define its deadline
$d_i=\max\{t\le r:A_i(t)\le1\}$. Give every remaining selected mode
the artificial deadline $r$. Sort these $r-2$ modes by nondecreasing
deadline, resolving ties arbitrarily, and assign them to the unit slots
with start times
\[
 0,1,\ldots,j-1,\qquad j+2,j+3,\ldots,r-1.
\]
We verify that each mode starts by its deadline. If the $\ell$th sorted
mode precedes $q$, it starts at $\ell-1$. A deadline smaller than this
would imply that all first $\ell$ modes have allocation strictly greater
than one at $\ell-1$, contradicting total allocation $\ell-1$. If it
follows $q$, it starts at $\ell+1$, with $\ell\ge j+1$. A deadline
smaller than $\ell+1$ would give allocation greater than one to the first
$\ell$ modes there. Also $d_q\le j+2\le\ell+1$, so $q$ has allocation
at least one there. Their combined allocation would exceed $\ell+1$.
An artificial deadline $r$ cannot be violated by a slot start.

For a singly used mode, $W_i-A_i\le1$ because its total service is one.
Before activation, $A_i-W_i\le1$ by the deadline property; after
activation, its largest positive discrepancy is at $r$ and is at most
$A_i(r)-1\le1$. For $q$, the allocation at its block start is at most
one and at its block end is at least one. Thus its positive discrepancy
at the start and its negative discrepancy at the end are at most one.
Its final positive discrepancy is $A_q(r)-2\le1$. Endpoint monotonicity
controls the intervening times. Every omitted mode has allocation at most
one throughout the prefix. The new prefix therefore has full error at
most one. Since its counts at $r$ are unchanged, attaching the original
suffix preserves every later discrepancy. The two $q$ slots are adjacent.
\end{proof}

\begin{proof}[Proof of Theorem~\ref{thm:heavy}]
Extend the input from $T$ to $(s+2)E$, when necessary, by assigning all
added activity to any mode. An originally heavy mode remains heavy.
Scale time and all allocations by $E$, giving the integer horizon
$M=s+2$. Lemmas~\ref{lem:prefix-repeat} and~\ref{lem:first-repeat}
produce an error-one word with an adjacent pair. It has at most $M-2=s$
switches. Scale back and restrict to $[0,T]$; neither operation increases
the switch count, and the error is at most $E$.
\end{proof}

The reordered word need not preserve the stronger floor/ceiling bounds
in~\eqref{eq:integral-prefix}; the conclusion needed is error at most one.
Nor does the proof prescribe the repeated mode: the first repeated mode
can differ from the heavy mode used in the flow objective.

\begin{theorem}[Exact one-sided reduction]\label{thm:one-sided-reduction}
For every $1\le k<n$,
\begin{equation}\label{eq:one-sided-reduction}
 F_{n,k-1}(T)=\max\left\{\frac{T}{k+1},G^-_{n,k}(T)\right\}.
\end{equation}
Repeated modes are allowed in both optimizations.
\end{theorem}
\begin{proof}
Set $E=\max\{T/(k+1),G^-_{n,k}(T)\}>0$. For an input with a mass
greater than $E$, Theorem~\ref{thm:heavy} applies with $s=k-1$. Otherwise
every mass is at most $E$. Choose a one-sided minimizing schedule, which
exists by Proposition~\ref{prop:compactness}; its negative discrepancy is
at most $G^-_{n,k}(T)\le E$. Its positive discrepancy is bounded by
$A_i(t)\le m_i\le E$, regardless of the schedule. This proves the upper
bound. Full error dominates one-sided error for every input and schedule,
so $F_{n,k-1}\ge G^-_{n,k}$. For the remaining lower bound, take $k+1$
successive distinct pure modes, each for time $T/(k+1)$. Every schedule
with at most $k$ blocks omits one of these modes and incurs its full
terminal mass as error. The assumption $k<n$ permits this input.
\end{proof}

\section{Analytic reach bounds with two and three distinct modes}\label{sec:reach-three}

Distinct-mode schedules are a useful construction for the one-sided
problem, although distinctness is never imposed on the competing optimum.
For a threshold $E>0$, a horizon $L>0$, and an unused mode $i$, define
\begin{equation}\label{eq:reach-definition}
 H_i(t)=t-A_i(t),\qquad
 \Phi_i(b)=\max\{t\in[0,L]:H_i(t)\le b+E\},\quad 0\le b\le L.
\end{equation}
The functions $H_i$ are continuous and nondecreasing, with
$\sum_iH_i(t)=(n-1)t$. The feasible set is closed, nonempty, and contains
$b$; hence $\Phi_i(b)\ge b$. Activating unused mode $i$ from $b$ to $t$
gives endpoint negative discrepancy $t-b-A_i(t)=H_i(t)-b$.
Its maximum negative discrepancy occurs at the end of that block.
Consequently a composition of reach maps along distinct modes constructs
a schedule with one-sided error at most $E$. Each map is nondecreasing in
$b$, so using the latest feasible endpoint at each step attains the greatest
reach for that word.

We will repeatedly use the exact boundary facts
\begin{equation}\label{eq:reach-boundary}
 \Phi_i(b)<L\ \Longrightarrow\ H_i(\Phi_i(b))=b+E,
 \qquad
 \Phi_i(b)<L\ \Longrightarrow\ H_i(L)>b+E.
\end{equation}
The first follows by continuity and maximality, and the second because
$L$ is then infeasible. Both remain valid with flat portions of $H_i$;
choosing an arbitrary threshold root instead of the latest endpoint would
not suffice for all comparisons below. Put
\begin{equation}\label{eq:B-def}
 r_n=\frac{n}{n-1},\qquad B_j(n,E)=nE(r_n^j-1),\qquad B_0(n,E)=0.
\end{equation}
When $n,E$ are fixed, write $B_j=B_j(n,E)$. These values satisfy
$B_j=r_n(B_{j-1}+E)$.

\begin{lemma}[Two distinct blocks]\label{lem:reach-two}
For $n\ge3$, two distinct modes suffice to reach $\min\{T,B_2\}$
with one-sided error at most $E$. A single mode suffices to reach
$\min\{T,B_1\}$ for every $n\ge2$.
\end{lemma}
\begin{proof}
Define reaches on the target horizon $L$. If a candidate already reaches
$L$, there is nothing to prove. Otherwise the first reaches
$R_i=\Phi_i(0)$ are uncapped, so $A_i(R_i)=R_i-E$.
For the single-block assertion, $x=\max_iR_i<L$ implies
$(n-1)x=\sum_iH_i(x)\ge nE$, so $x\ge B_1$, a contradiction
when $L\le B_1$.

For two blocks, assume every ordered distinct pair has reach below $L$,
and let $M=\max_{p\ne q}\Phi_q(R_p)<L$. Let $x\ge y$ be the
two largest first reaches and choose $p$ attaining $x$. At time $y$,
$A_p(y)\le x-E$ and $A_i(y)\le y-E$ for every $i\ne p$.
Summing gives $x+(n-2)y\ge nE$. The identity
\begin{equation}\label{eq:first-reach-pair}
 (n-1)x+y=\frac{n}{n-1}[x+(n-2)y]
 +\frac{n^2-3n+1}{n-1}(x-y)
 \ge\frac{n^2E}{n-1}
\end{equation}
uses $n\ge3$. For each $i$, let $x_i=\max_{j\ne i}R_j$.
Every pair ending in $i$ reaches no further than $M$, so
\eqref{eq:reach-boundary} and monotonicity give
$A_i(M)\le M-E-x_i$. Sum over $i$ and note
$\sum_i x_i=(n-1)x+y$. It follows that
\begin{equation}\label{eq:global-pair-bound}
 (n-1)M\ge nE+(n-1)x+y,
 \qquad M\ge\frac{nE(2n-1)}{(n-1)^2}=B_2.
\end{equation}
This contradicts $M<L\le B_2$ and proves the claim.
\end{proof}

\begin{theorem}[Three distinct blocks]\label{thm:reach-three}
For $n\ge4$, three distinct modes suffice to reach $\min\{T,B_3\}$
with one-sided error at most $E$.
\end{theorem}
\begin{proof}
Take $L=\min\{T,B_3\}$ and suppose no sequence of at most three
distinct modes reaches $L$. In particular every first and pair reach is
uncapped. Retain $R_i,M,x,y$ from Lemma~\ref{lem:reach-two}, and let
$z$ be the third largest first reach. At time $z$, the top two modes
have allocation at most $x-E$ and $y-E$, and the other modes have
allocation at most $z-E$. Thus $x+y+(n-3)z\ge nE$.
Writing $u=x+y\ge2z$ gives
\begin{equation}\label{eq:first-reach-triple}
 \begin{split}
 (n-1)u+2z
 &=\frac{2n}{n-1}[u+(n-3)z]
 +\frac{n^2-4n+1}{n-1}(u-2z)\\
 &\ge\frac{2n^2E}{n-1},
 \end{split}
\end{equation}
where the last coefficient is positive for $n\ge4$.

Let $M_i$ be the largest pair reach excluding mode $i$, and let
$x_i\ge y_i$ be the largest two first reaches after that exclusion.
For each $a\ne i$, all pairs available without $i$ and ending in $a$
give
\[
 A_a(M_i)\le M_i-E-\max_{b\notin\{i,a\}}R_b.
\]
The maxima on the right sum over $a\ne i$ to $(n-2)x_i+y_i$.
Summing and using $\sum_a A_a(M_i)=M_i$ yields
\[
 A_i(M_i)\ge(n-1)E+(n-2)x_i+y_i-(n-2)M_i.
\]
On the other hand, $M_i\le M$ and
$A_i(M)\le M-E-x_i$. Monotonicity therefore implies
\begin{equation}\label{eq:excluded-pair-ineq}
 (n-2)M_i\ge nE+(n-1)x_i+y_i-M.
\end{equation}

Choose a pair $(p,q)$ attaining $M$. The same pair is available in every
exclusion except possibly $p,q$, so $M_i=M$ outside this pair. Also
$x_p+x_q\ge x+y$ and $y_p,y_q\ge z$. Summing
\eqref{eq:excluded-pair-ineq} for $p,q$ gives
\[
 \sum_i M_i\ge
 \left(n-2-\frac2{n-2}\right)M
 +\frac{2nE+(n-1)(x+y)+2z}{n-2}.
\]
The coefficient of $M$ is positive for $n\ge4$. Substituting
$M\ge B_2$ from~\eqref{eq:global-pair-bound} and
\eqref{eq:first-reach-triple} proves
\begin{equation}\label{eq:aggregate-pair}
 \sum_iM_i\ge nB_2.
\end{equation}
For completeness, the substituted expression equals $nB_2$ because
$(n-1)B_2=nE+n^2E/(n-1)$.

Appending mode $i$ to a maximizing pair excluding it fails to reach $L$
by assumption. The second part of~\eqref{eq:reach-boundary} gives
$H_i(L)>M_i+E$. Sum these strict inequalities and use
\eqref{eq:aggregate-pair}:
\[
 (n-1)L>\sum_iM_i+nE\ge nB_2+nE=(n-1)B_3.
\]
This contradicts $L\le B_3$.
\end{proof}

This proof is constructive: compute the first reaches and the ordered-pair
reach matrix, then append each possible final mode to a maximizing pair
excluding it. A global maximizing pair supplies all but at most two of
these excluded maxima; the remaining two require additional scans of the
matrix. The construction therefore uses $n(n-1)$ pair reach evaluations,
followed by $n$ final reach evaluations and $O(n^2)$ comparisons. It
replaces an earlier proof for $n=4$ based on thirty event-order linear
programs. Those exact certificates are included as an independent
special-case check in the supplement, but are not premises of
Theorem~\ref{thm:reach-three}.
